\documentclass[]{scrartcl}

\title{MIPS Home Computer System}
\subtitle{System Design Reference}
\author{
    \large Student\\
    \large Mostafa Abd El-Aziz\\
    \and
    \large Supervisor\\
    \large Prof. Dr. Layla Abou Hadid\\
}
\date{\vspace{-5ex}}

% font
\usepackage[T1]{fontenc}
\usepackage{lmodern}
\renewcommand*\familydefault{clms}
\usepackage{sfmath}
\usepackage{microtype}
\usepackage[utf8]{inputenc}
\usepackage{cfr-lm} % font package

% listings
\usepackage{listings}
\lstset{language=java}

% margins
\usepackage[margin=3cm,a4paper]{geometry}

% figures and tables
\usepackage{graphicx}

% table of contents
\usepackage{tocloft}
\renewcommand{\cftsecleader}{\cftdotfill{\cftdotsep}} % for chapters
\renewcommand{\cftpartleader}{\cftdotfill{\cftdotsep}} % for parts

% section style
\usepackage{sectsty}
\sectionfont{\LARGE}

% page orientation
\usepackage{lscape}

%------------------------------------------------------------------------------%

\begin{document}

% university logo
\titlehead{

    \small Alexandria University\\
    \small Faculty of Engineering\\
    \small Computer and Systems Engineering Department\\
    \small CS 401 - Graduation Project

    \vspace{-2.2cm}
    \hfill
    \includegraphics[scale=0.6]{logo.png}
}

% document title
\maketitle

% table of contents
\renewcommand{\contentsname}{Table of Contents}
\tableofcontents

%------------------------------------------------------------------------------%

\section{Introduction}

The project is implemented using Nexys-2 FPGA board. The board contains
a Spartan-3e FPGA chip, RAM chip, ROM chip, 8 input switches, 4 push buttons, 4
7-segment LEDs, VGA port, RS-232 port, PS/2 port, 32 General Purpose I/O pins,
50MHz crystal oscillator, JTAG, USB (for interfacing with a PC), and other
relevant components.\\

We can divide the design of the system into two main parts:

\begin{itemize}
\item \textbf{Solid part}: the solid components on Nexys-2 board that are
      used as functional part of the computer system. The components are:
      \begin{itemize}
      \item \underline{On-board ROM chip}: stores computer firmware.
      \item \underline{On-board RAM chip}: mainly used to store software that
            is being executed by CPU (except firmware, which is
            executed directly from ROM) and temporary data.
      \item \underline{VGA port}: a VGA-compatible CRT screen is connected to
            the VGA port to display information.
      \item \underline{PS/2 port}: a PS/2 keyboard is connected to PS/2 port
            to feed the machine with data.
      \item \underline{50MHz crystal}: responsible for synchronization of
            the system.
      \end{itemize}
\item \textbf{Volatile part}: implemented in VHDL and compiled using Xilinx
      SDK. The compiled design is uploaded to Spartan-3e FPGA chip through
      USB and JTAG. The logic uploaded to FPGA communicate with the
      outer world (components on Nexys II boards) through FPGA chip pins
      that are connected to this components through wires. This part
      consists of the following HDL modules:
      \begin{itemize}
      \item \underline{Central Processing Unit}:
            implements MIPS I instruction set.
      \item \underline{Memory Interface Unit}:
            interfaces on-board RAM and ROM.
      \item \underline{VGA Controller}: interface on-board VGA port.
      \item \underline{PS/2 Controller}:
            interface on-board PS/2 port (not implemented).
      \item \underline{Bus controller}:
      \     logic of communication between CPU and other components.
      \end{itemize}
\end{itemize}

When the computer boots up, CPU starts executing firmware instructions
that are stored on the ROM chip on Nexys-II board. Software can access
I/O devices (like VGA controller) using memory instructions (load and
store). The CPU doesn't differentiate between I/O accesses and memory accesses
(there is no isolated I/O space). Consequently, I/O devices share the
same bus with RAM.\\

Memory locations that are reserved for communication with I/O devices are
called memory-mapped I/O registers. The following table shows the memory
layout and memory-mapped I/O locations:\\

\begin{figure}[h]
\begin{center}
  \def\arraystretch{1.3}
  \begin{tabular}{ |@{\hspace{1em}} c @{\hspace{1em}} | c | c | }
    \hline
    \textbf{Address Region} & \textbf{Region Size} & \textbf{Device} \\
    \hline
    0x00000000-0x0000FFFF & 64KB & Lowest 64KB of on-board ROM \\
    \hline
    0x00010000-0x00017FFF & 32KB & Lowest 32KB of on-board RAM \\
    \hline
    0x00018000-0x0001FFFF & 32KB & VGA frame buffer \\
    \hline
    0xFFF00000-0xFFFFFFFF & 1MB & Keyboard I/O \\
    \hline
  \end{tabular}
  \caption{Memory Address Space layout}
\end{center}
\end{figure}

When CPU needs to fetch an instruction or a data word from memory, the
location of this instruction/data is put onto the address bus. Logic
on the FPGA decides which component shall reply with data based on the table
shown above. This logic is part of top-level design of the machine. Sheet 1
(appended to this documents) shows the high-level components of the system.

%------------------------------------------------------------------------------%

\section{Central Processing Unit}

The CPU module implements MIPS I instruction set architecture along with
an instruction cache and a data cache. It is planned to implement a TLB in the
future. Logically, the module consists of two submodules: the core pipeline and
the cache system (as shown in sheet 2).

\subsection{The Core Pipeline}

This module contains the instruction path and data path which are pipelined
into 5 stages: IF, ID, EX, MEM, and WB. IF is responsible for interfacing
instruction cache, while ID contains register file, hazard detection unit,
and instruction decoding logic. EX contains ALU, MEM interface data cache,
and finally, WB interfaces register file for write.\\

The pipeline is an extended version of the simple 5-stage pipeline of MIPS
subset. Our proposed design completely implements MIPS I ISA as we said
before. The pipeline is written in VHDL as a single module that is made
to interface the cache directly. Sheet 3 shows the design of the pipeline.

\subsection{Cache System}

The cache module consists of an instruction cache, a data cache, a buffer,
and a sequential control unit that detects cache misses and fills in the buffer
with memory orders. As long as the buffer is not empty, the pipeline is
stalled and memory cycles are issued to execute the entries of the buffer.\\

Both instruction and data caches are write-through caches, so any store
instruction will eventually stall the pipeline until memory write
cycle is done (and the buffer becomes empty). The cache array is
implemented using block memory on the FPGA chip.\\

Cache module is a multiplexing unit. Instruction cache hit and data cache hit
can be processed in the same time, but if instruction and data cache misses
occur in the same time, instruction cache miss will be processed first, then
data cache miss. This limitation is due to the single-bus nature of Von Neumann
architecture (Von Neumann Bottleneck).

\section{Memory Interface}

The memory interface module is responsible for controlling the bus between
FPGA chip and on-board ROM and RAM chips. Each chip has 16MB memory, arranged
in 8M 16-bit words. This results in inconsistency between word size of both ICs
and the word size of CPU data bus.\\

To deal with this problem, the memory interface module translate 32-bit
read/write cycles into two 16-bit read/write orders that are processed
in sequence. 8- and 16-bit cycles don't need to be divided into two cycles,
so they can finish up faster.

\section{VGA Controller}

The VGA module is designed to have the same software interface of IBM PC VGA in
text-mode. In PC, VGA provides a buffer (starting at 0xB8000) that can be
accessed directly by software to draw characters on screen (text mode). In mode
03h, screen has 25 rows and 80 columns, resulting in 2000 characters (4000
bytes).\\

Our implementation provides this interface by implementing several VGA RAM
modules that are mapped to specific memory locations (0x18000-0x1FFFF). The
VGA RAM is interfaced to system bus through graphics controller. VGA contains
logic (the sequencer and CRT controller) that continuously scans VGA RAM and
generates RGB signals to refresh CRT screen view. Nexys II provides an On-chip
DAC (a set of resistors) to convert RGB signals into analog form.

\end{document}
